\graphicspath{{images/}}
\renewcommand{\theequation}{\arabic{chapter}.\arabic{equation}}

\chapter{Marco metodológico}

Este capítulo presenta de forma concisa los pasos técnicos y los
requisitos necesarios para levantar un entorno de simulación de
tráfico con SUMO y preparar el escenario que servirá de base para
entrenar un agente de aprendizaje por refuerzo.  El objetivo es que
cualquier lector pueda instalar SUMO, importar una intersección real
mediante datos de OpenStreetMap (OSM), generar la red y las rutas
vehiculares, y ejecutar la simulación desde la línea de comandos.  Se
incluye además una descripción breve de las funciones de las unidades
U1 (preprocesamiento) y U2 (entrenamiento) y los requisitos de
hardware mínimos recomendados.

\section{Alcance y unidades}

\noindent\textbf{U1: Preprocesamiento.} Comprende todas las tareas
relacionadas con la construcción del escenario.  Esto incluye la
instalación y configuración de SUMO, la descarga y conversión de un
recorte de OSM en una red SUMO, la generación de rutas básicas de
vehículos y la creación de archivos de configuración reproducibles
(\texttt{.net.xml} para la red, \texttt{.rou.xml} para la demanda y
\texttt{.sumocfg} para el escenario).  El resultado de U1 es un
escenario listo para simular.

\noindent\textbf{U2: Entrenamiento.} Toma como entrada el escenario de
U1 y ejecuta SUMO a través de la interfaz TraCI para que un agente de
aprendizaje por refuerzo observe el estado del tráfico y modifique
las fases de los semáforos.  En U2 no se modifican la red ni las
rutas: el enfoque se centra en el algoritmo de control.

\section{Requisitos de hardware}

Para ejecutar U1 y U2 se recomiendan los recursos de la
Tabla~\ref{tab:hw}.  U1 puede ejecutarse en un equipo doméstico
modesto, mientras que U2 se beneficia de una GPU si se entrenan
modelos complejos.

\begin{table}[H]
\centering
\begin{tabular}{p{3.5cm}p{5cm}p{5cm}}
\hline
 & \textbf{Mínimo (U1)} & \textbf{Recomendado (U1+U2)} \\
\hline
CPU & 4 hilos reales & 8--12 hilos \\
RAM & 8~GB & 16--32~GB \\
GPU & no requerida & $\geq$6~GB VRAM \\
Almacenamiento & 10~GB SSD & 20--50~GB SSD \\
Sistema operativo & Windows 10/11 o Ubuntu 20.04+ & Windows 10/11 o Ubuntu 20.04+ \\
\hline
\end{tabular}
\caption{Recursos recomendados para construir y ejecutar el escenario.}
\label{tab:hw}
\end{table}

\section{Instalación y configuración de SUMO}

\begin{enumerate}
  \item \textbf{Descargar e instalar SUMO.} Obtenga la versión estable
    de SUMO desde el sitio oficial y siga las instrucciones de
    instalación para su sistema operativo.  SUMO incluye varios
    ejecutables de consola, como \emph{sumo} (simulación),
    \emph{sumo-gui} (versión con interfaz gráfica) y
    \emph{netconvert} (importación de redes).
    Instale también Python~3, necesario para ejecutar los
    scripts auxiliares.

  \item \textbf{Configurar el entorno.} Defina la variable de
    entorno \verb|SUMO_HOME| apuntando al directorio donde se instaló
    SUMO y añada la carpeta \verb|bin| al \verb|PATH| del sistema.
    Esto permite invocar los programas de SUMO desde cualquier
    terminal.

  \item \textbf{Verificar la instalación.} Abra una terminal y ejecute
    los siguientes comandos.  Deben devolver la versión instalada de
    SUMO para confirmar que el entorno está listo:

    \begin{verbatim}
sumo -V
sumo-gui -V
    \end{verbatim}
\end{enumerate}

%----------------------------------------------------------------------
\section{Limpieza y configuración del escenario OSM $\rightarrow$ SUMO}

Tras importar la red desde OSM, es necesario revisar y ajustar el
escenario para que represente fielmente la intersección de estudio y
resulte adecuado para el entrenamiento del agente de control. Este
proceso consiste en eliminar elementos irrelevantes, corregir
inconsistencias topológicas y configurar adecuadamente los semáforos y
la demanda. A continuación se describe un procedimiento general:

\subsection*{Eliminación y simplificación de la red}

\begin{enumerate}
  \item \textbf{Abrir la red en un editor.} Utilice \verb|sumo-gui| o
    \texttt{netedit}, las herramientas de SUMO destinadas a la
    inspección y edición de redes. Desde la interfaz gráfica se puede
    seleccionar manualmente cada calle, carril e intersección para
    examinar su geometría y atributos.

  \item \textbf{Eliminar calles irrelevantes.} Desactive o borre los
    segmentos de vía que no formen parte del cruce objeto de estudio.
    En \texttt{netedit} esto puede realizarse en modo ``Edge''
    seleccionando los elementos y pulsando la tecla Suprimir. También
    se pueden modificar jerarquías de carreteras para mantener
    únicamente arterias principales y desechar caminos secundarios.

  \item \textbf{Ajustar sentidos de circulación.} Verifique que las
    direcciones de cada carril coinciden con la realidad (por ejemplo,
    si un carril debería ser unidireccional). En \texttt{netedit}
    es posible invertir la dirección de un carril o fusionar dos
    carriles paralelos empleando la herramienta de edición de lanes.

  \item \textbf{Simplificar nodos complejos.} Algunas importaciones
    desde OSM generan nodos excesivamente subdivididos. Estos pueden
    fusionarse mediante la opción de simplificación de intersecciones
    en \texttt{netconvert} o mediante edición manual. Una red más
    simple facilita el entrenamiento y reduce tiempos de simulación.
\end{enumerate}

Tras completar estos pasos, la red se limpia y simplifica. La
Figura~\ref{fig:netedit-before-after} muestra un ejemplo antes y
después de la limpieza realizada en \texttt{netedit} para el cruce
estudiado. Se observa cómo se eliminan calles secundarias y se ajusta
la orientación de los carriles.

\begin{figure}[H]
  \centering
  \subfigure[Vista inicial en \texttt{netedit} con la intersección
    importada; hay calles secundarias y direcciones de carril
    incorrectas.]{\includegraphics[width=0.48\textwidth]{NETEDIT.PNG}}
  \hfill
  \subfigure[Misma intersección tras la limpieza: calles irrelevantes
    eliminadas y sentidos de los carriles corregidos.]{\includegraphics[width=0.48\textwidth]{conexiones.PNG}}
  \caption{Antes y después de la limpieza de la red en \texttt{netedit}.}
  \label{fig:netedit-before-after}
\end{figure}

\subsection*{Configuración de semáforos y demanda}

\begin{enumerate}
  \item \textbf{Revisar semáforos importados.} Si al generar el
    escenario se activó la opción ``Traffic lights'', los semáforos
    aparecerán como nodos \verb|<tlLogic>| en el archivo
    \texttt{.net.xml}. Utilice \emph{sumo-gui} o \texttt{netedit}
    para inspeccionar las fases (colores) y tiempos asignados
    automáticamente. Ajuste los valores para que respeten las
    normativas locales: por ejemplo, establezca un mínimo de cuatro
    segundos de luz ámbar y defina duraciones de luz verde
    proporcionales al flujo de cada aproximación.

  \item \textbf{Crear o modificar programas de semáforos.} Para
    controlar manualmente los semáforos se puede definir un archivo
    adicional \texttt{.add.xml}. Un ejemplo sencillo de un programa
    estático para una intersección de dos fases es:

\begin{verbatim}
<additional>
  <tlLogic id="tl_cruce" type="static" programID="base" offset="0">
    <!-- Fase 0: verde en NS, rojo en EW -->
    <phase duration="30" state="GrGr"/>
    <!-- Fase 1: ámbar en todas las aproximaciones -->
    <phase duration="4"  state="yryr"/>
    <!-- Fase 2: verde en EW, rojo en NS -->
    <phase duration="30" state="rGrG"/>
    <!-- Fase 3: ámbar en todas las aproximaciones -->
    <phase duration="4"  state="ryry"/>
  </tlLogic>
</additional>
\end{verbatim}

    Este archivo debe referenciarse en el fichero \texttt{.sumocfg}
    dentro del apartado \verb|<additional-files>| para que SUMO lo
    cargue junto con la red y las rutas. Si se decide mantener los
    semáforos generados automáticamente, conviene al menos fijar sus
    tiempos mínimos y comprobar que las fases son coherentes.

  \item \textbf{Configurar la demanda.} La generación de rutas con
    \verb|randomTrips.py| produce una distribución uniforme de viajes
    sobre la red. Para configurar tasas de llegada realistas
    conviene ajustar el parámetro \verb|-p| (periodo medio entre
    vehículos) y usar distintas semillas con \verb|--seed| para
    reproducibilidad. Si se dispone de datos de tráfico, se pueden
    definir flujos en un archivo \texttt{.rou.xml} indicando el
    intervalo de tiempo, la tasa de arribo y el tipo de vehículo. Un
    ejemplo básico de flujo es:

\begin{verbatim}
<routes>
  <vType id="car" accel="2.6" decel="4.5" length="5" maxSpeed="13.9"/>
  <flow id="f1" begin="0" end="900" vehsPerHour="600" type="car"
        from="edgeWest" to="edgeEast"/>
</routes>
\end{verbatim}

    Estas definiciones permiten controlar la intensidad y dirección del
    tráfico, así como simular distintos tipos de vehículos (coches,
    autobuses, motocicletas, etc.). Tras ajustar la demanda, use
    \texttt{duarouter} para convertir los flujos en rutas específicas
    si fuera necesario.

  \item \textbf{Actualizar los archivos de configuración.} Una vez
    depurada la red y configurados los semáforos y la demanda,
    asegúrese de que el archivo de configuración \texttt{.sumocfg}
    referencia correctamente todos los ficheros relevantes:
    \texttt{.net.xml}, \texttt{.rou.xml} y, en su caso,
    \texttt{.add.xml}. Ejecute una simulación breve para verificar
    que el escenario se carga sin mensajes de error.
\end{enumerate}

\subsection*{Herramientas útiles}

Para tareas de limpieza y configuración puede emplear las siguientes
herramientas que acompañan a SUMO:

\begin{itemize}
  \item \textbf{netedit:} editor gráfico que permite crear, modificar
    y depurar redes, así como editar semáforos, conexiones y
    restricciones de circulación. Resulta idóneo para eliminar
    segmentos irrelevantes y ajustar sentidos de vía.
  \item \textbf{sumo-gui:} interfaz gráfica de simulación que permite
    visualizar el comportamiento de los vehículos y los semáforos en
    tiempo real. Es útil para validar la red depurada y detectar
    cuellos de botella o comportamientos anómalos.
  \item \textbf{duarouter:} herramienta para convertir definiciones de
    flujos (\verb|<flow>|) en rutas concretas, realizando el cálculo
    del recorrido óptimo entre origen y destino. Es necesaria cuando
    se trabaja con demanda no uniforme basada en datos observados.
\end{itemize}

\section{Importación de un cruce desde OSM}

Una vez instalado SUMO, el siguiente paso consiste en importar la
intersección o corredor de interés.  Existen dos alternativas.

\subsection{Asistente OSM (\emph{osmWebWizard})}

\begin{enumerate}
  \item Ejecute el script \verb|osmWebWizard.py| desde la carpeta de
    herramientas de SUMO.  El script abre un navegador donde se puede
    navegar sobre el mapa de OSM, seleccionar la zona deseada y
    generar el escenario automáticamente.  El uso básico es:
    \begin{verbatim}
python tools/osmWebWizard.py
    \end{verbatim}

  \item Seleccione ``Lima, Peru'' en el buscador y ajuste el
    \emph{bounding box} sobre la intersección de interés.  Active la
    opción “Traffic lights” si desea que se importen los semáforos y
    pulse “Generate Scenario”.  El asistente creará la red
    \texttt{.net.xml}, una demanda aleatoria y un fichero
    \texttt{.sumocfg}; además iniciará \emph{sumo-gui} con el
    escenario resultante.  Los archivos se guardan en un
    subdirectorio con marca temporal.

  \item Revise la red en \emph{sumo-gui} y, si detecta calles
    desconectadas o giros imposibles, ajuste el recorte y repita la
    generación hasta obtener una red consistente.
\end{enumerate}

Para ilustrar este proceso, la Figura~\ref{fig:import-osm-sumo} muestra
la selección del área en OSM y la red generada en SUMO. A la
izquierda se observa el recuadro delimitado sobre el mapa en el asistente
\emph{osmWebWizard}, y a la derecha la red importada tal como aparece en
\emph{sumo-gui}.

\begin{figure}[H]
  \centering
  \subfigure[Selección del \textit{bounding box} en OSM mediante el
    asistente \emph{osmWebWizard}.]{\includegraphics[width=0.48\textwidth]{OSM.PNG}}
  \hfill
  \subfigure[Red del cruce generada automáticamente en \emph{sumo-gui}.]{\includegraphics[width=0.48\textwidth]{SUMO.PNG}}
  \caption{Importación de un cruce desde OSM a SUMO.}
  \label{fig:import-osm-sumo}
\end{figure}

\section{Generación y ejecución del escenario}

\subsection{Generar rutas de vehículos}

Con la red validada, se requieren rutas de vehículos para simular el
tráfico:

\begin{enumerate}
  \item Para un escenario preliminar se puede recurrir a
    \verb|randomTrips.py|.  Este script selecciona aleatoriamente
    arcos de la red como origen y destino, distribuyendo los vehículos
    entre un tiempo de inicio \verb|-b| y uno de fin \verb|-e|.
    Se controla la intensidad del flujo con la opción \verb|-p| y se
    garantiza reproducibilidad fijando una semilla con \verb|--seed|.
    Un ejemplo básico que genera viajes durante 15 minutos con un
    vehículo por segundo es:

    \begin{verbatim}
python tools/randomTrips.py -n lima.net.xml -b 0 -e 900 -p 1.0 \
    --route-file rutas.rou.xml
    \end{verbatim}

  \item Si se dispone de datos observados (aforos o tiempos de viaje),
    se pueden definir flujos origen–destino y enrutar con
    \verb|duarouter|.  Este proceso requiere calibrar tasas de arribo
    y cae fuera del alcance de este procedimiento básico.
\end{enumerate}

\subsection{Crear el archivo de configuración}

Agrupe los archivos de red (\texttt{.net.xml}), rutas
(\texttt{.rou.xml}) y, si procede, el programa de semáforos
(\texttt{.add.xml}) en un único archivo de configuración
\texttt{.sumocfg}.  Este fichero referencia todos los elementos de
entrada y define los intervalos de simulación.

\subsection{Ejecución de SUMO}

\begin{enumerate}
  \item Para simular de forma interactiva, ejecute SUMO con
    interfaz gráfica:
    \begin{verbatim}
sumo-gui -c escenario.sumocfg
    \end{verbatim}

  \item Para simulaciones por lotes sin interfaz gráfica, utilice el
    ejecutable \verb|sumo|:
    \begin{verbatim}
sumo -c escenario.sumocfg
    \end{verbatim}
    Los ejecutables de SUMO son aplicaciones de consola; las opciones
    \verb|--net| y \verb|--route-files| permiten especificar
    directamente los archivos de red y rutas si no se usa un
    \texttt{.sumocfg}.
\end{enumerate}

\section{Escenarios sintéticos de prueba}
\label{sec:escenarios}

Mientras se completa la importación del cruce real de Lima, los
módulos U2 y U3 se desarrollaron y validaron sobre cuatro escenarios
sintéticos, construidos con \texttt{netconvert} a partir de archivos de
nodos y vías escritos a mano, del más simple al más complejo
(Tabla~\ref{tab:escenarios}). Todos comparten el mismo tipo de
vehículo (turismo de 5~m, aceleración 2.6~m/s$^2$, velocidad máxima
50~km/h, imperfección del conductor $\sigma = 0.5$) y el mismo patrón
de demanda: una hora dividida en dos periodos de 30 minutos con la
punta invertida, de modo que un programa de tiempo fijo no pueda ser
óptimo en ambos a la vez. Los flujos se definen con
\verb|<flow ... vehsPerHour>| por movimiento y periodo, lo que produce
llegadas equiespaciadas; el ámbar es de 3~s en todos los casos (valor
por defecto de \texttt{netconvert}).

\begin{table}[H]
\centering
\small
\setlength{\tabcolsep}{3pt}
\begin{tabular}{p{1.3cm}p{3.4cm}p{1.4cm}p{3.4cm}p{1.2cm}p{2.1cm}}
\toprule
Escenario & Geometría & Fases verdes & Demanda del periodo punta (veh/h) & Vehículos en la hora & Programa fijo de referencia \\
\midrule
cruce & cruce N--S / E--O, 1 carril por sentido & 2 & eje principal 650 por sentido; transversal 150; giros 40 a 80 & 1844 & verde 25~s, ciclo 56~s \\
cruce2 & 2 carriles por sentido más carril exclusivo de giro a la izquierda & 4 (2 principales y 2 de giro protegido) & recto 420 por sentido; giro izquierdo 160; transversal 200 & 1884 & verde 20~s, giro 6~s, ciclo 64~s \\
corredor & dos cruces separados 300~m sobre una avenida E--O, 1 carril por sentido & 2 por cruce & avenida 600 en el sentido punta y 250 en el contrario; transversales 180 & 1810 & verde 25~s; desfase de 25~s en el segundo cruce (onda verde) \\
red & rotonda sin semáforo más tres cruces separados 300~m sobre una avenida de 2 carriles por sentido & 2 por cruce & avenida 650 y 250; transversales 150; rotonda 80 a 150 & 2152 & verde 20~s; desfases 0, 23 y 0~s; ciclo 46~s \\
\bottomrule
\end{tabular}
\caption{Escenarios sintéticos de prueba. La demanda del segundo periodo es la del primero con la punta invertida.}
\label{tab:escenarios}
\end{table}

El programa fijo de cada escenario no es un baseline débil. El verde y,
en los corredores, el desfase entre semáforos se ajustaron primero por
barrido de la espera detenida y después se comprobaron con un barrido
por retraso promedio sobre las semillas 1001 a 1003
(\texttt{barrer\_fijo.py}, Anexo~B), cuya malla incluye la
configuración usada: en cruce2, corredor y red esta resultó ser la de
menor retraso de la malla, y en cruce queda 0.2~s por encima del verde
de 24~s, dentro de la variación entre semillas. En red el desfase de
23~s coincide con el tiempo de viaje entre cruces. Como segunda referencia se usa el control actuado
nativo de SUMO, con las mismas fases y los mismos límites de verde
mínimo y máximo que el agente (Sección~\ref{sec:plan}).

\section{Diagrama de flujo del módulo U1}

La Figura~\ref{fig:flow-u1} resume el procedimiento completo del módulo U1,
desde la selección del cruce en Lima hasta la obtención del escenario SUMO
listo para ser utilizado en el entrenamiento del agente de RL (U2).

\begin{figure}[H]
    \centering
    \includegraphics[width=\textwidth]{U1.png}
    \caption{Diagrama de flujo del módulo U1: preprocesamiento y generación del
    escenario SUMO.}
    \label{fig:flow-u1}
\end{figure}

\section{Resumen de las unidades}

\begin{description}
  \item[U1:] Instalar y configurar SUMO; importar la intersección
    mediante el asistente OSM o descarga manual; limpiar y validar
    la red; generar la demanda vehicular y el archivo de configuración
    \texttt{.sumocfg}.  El resultado es un escenario reproducible.
  \item[U2:] Ejecutar el escenario con TraCI y entrenar un agente de
    aprendizaje por refuerzo que ajuste las fases semafóricas.  U2
    depende exclusivamente de los archivos generados en U1 y no
    modifica la red ni las rutas.
  \item[U3:] Medir el desempeño de cada control con las mismas salidas
    de SUMO (viajes por vehículo, colas por TraCI, emisiones) y
    compararlo con el protocolo del plan de pruebas
    (Secciones~\ref{sec:u3} y~\ref{sec:plan}).
  \item[U4:] Coordinar los agentes de semáforos vecinos. Pendiente: en
    los escenarios de varios cruces los agentes son, por ahora,
    independientes.
\end{description}

%----------------------------------------------------------------------
\section{Desarrollo del módulo U2: Aprendizaje por Refuerzo para el control de semáforos}

El segundo módulo de la metodología aborda el diseño e
implementación de un controlador inteligente que modifique las fases
de los semáforos en función del estado del tráfico. Para ello se
recurre a técnicas de \emph{Aprendizaje por Refuerzo} (\emph{Reinforcement
Learning}, RL), en las que un agente aprende a tomar decisiones
mediante la interacción con un entorno simulado. Este apartado
describe la formulación del problema, la interfaz con SUMO y los
algoritmos de RL más habituales para el control de intersecciones.

\subsection*{Formulación del problema}

En RL se define un \emph{entorno} compuesto por un conjunto de
\emph{estados}, un conjunto de \emph{acciones} y una función de
\emph{recompensa}. La formulación implementada en esta tesis para un
cruce con cuatro aproximaciones es la siguiente:

\begin{itemize}
  \item \textbf{Estado:} el número de vehículos detenidos en cada
    aproximación, leído con \verb|edge.getLastStepHaltingNumber| y
    discretizado en cuatro niveles (0; 1 a 3; 4 a 8; más de 8), más el
    índice de la fase verde activa. El estado es $s = (c_N, c_S, c_E,
    c_O, f)$ con $c_i \in \{0,1,2,3\}$. Con dos fases verdes el espacio
    tiene $4^4 \times 2 = 512$ estados y con cuatro fases, $1024$. Esta
    discretización es la que hace viable una tabla Q y, a la vez, la
    que limita al método tabular cuando crece el número de fases, como
    se muestra en el Capítulo~\ref{cap:resultados}.

  \item \textbf{Acción:} dos acciones. $a = 0$ mantiene la fase verde
    actual y $a = 1$ pasa a la siguiente fase verde del programa,
    insertando automáticamente el ámbar del programa (3~s). El agente
    no elige una fase arbitraria ni una duración: decide únicamente
    cuándo termina el verde en curso.

  \item \textbf{Recompensa:} negativa y proporcional a la congestión en
    el instante de la decisión,
    \begin{equation}
      r = -\Bigl(w_1 \sum_i c_i + w_2 \sum_i e_i\Bigr), \qquad
      w_1 = 1.0,\; w_2 = 0.01,
      \label{eq:recompensa}
    \end{equation}
    donde $c_i$ es el número de vehículos detenidos en la aproximación
    $i$ (sin discretizar) y $e_i$ es la suma de los tiempos de espera de
    los vehículos que en ese momento están en la aproximación $i$
    (\verb|edge.getWaitingTime|). Con estos pesos domina el término de
    colas; el de espera desempata entre estados con la misma cola. La
    recompensa mide el nivel de congestión, no su incremento entre
    decisiones.

  \item \textbf{Restricciones operativas:} el verde dura al menos 10~s
    antes de que el agente pueda cambiarlo (5~s en las fases de giro
    protegido, identificadas por tener a lo más dos movimientos en
    verde) y como máximo 60~s, tras los cuales el cambio se fuerza. El
    agente decide cada 5~s.
\end{itemize}

Con estos elementos se modela el control del semáforo como un
\emph{problema de decisión de Markov} (MDP). El agente observa el
estado, aplica una acción, avanza la simulación y recibe una
recompensa; a partir de la secuencia de experiencias va aprendiendo
una política de control. En los escenarios con varios cruces se
instancia un agente por semáforo, cada uno con su propia tabla Q, su
estado local y su recompensa local; los agentes no intercambian
información (eso es lo que añade el módulo U4).

\subsection*{Interfaz con SUMO y ciclo de entrenamiento}

Para implementar el bucle de RL se emplea la interfaz TraCI,
proporcionada por SUMO, que permite controlar la simulación desde
Python. El ciclo implementado consta de los siguientes pasos:

\begin{enumerate}
  \item \textbf{Inicio del episodio:} se lanza SUMO con
    \verb|traci.start| sobre el escenario de U1, con una semilla
    explícita (\verb|--seed|) y la salida por vehículo \texttt{tripinfo}
    activada. Se lee el programa del semáforo para identificar sus
    fases verdes y el ámbar que sigue a cada una, y se fija la primera
    fase verde con duración indefinida: a partir de ahí el agente
    decide cuándo cambia.
  \item \textbf{Decisión cada 5~s:} se lee el estado (colas
    discretizadas y fase activa) y se elige la acción con una política
    $\varepsilon$-greedy: con probabilidad $\varepsilon$ una acción al
    azar y, si no, la de mayor valor Q, desempatando al azar cuando
    ambas valen lo mismo. Si la acción es cambiar y el verde lleva al
    menos el mínimo, se activa el ámbar con \verb|trafficlight.setPhase|
    y \verb|setPhaseDuration|, se avanza la simulación esos segundos y
    se entra al siguiente verde. Luego se avanzan 5~s con
    \verb|simulationStep|.
  \item \textbf{Recompensa y actualización:} con el estado resultante
    se calcula la recompensa de la ecuación~\ref{eq:recompensa} y, en
    entrenamiento, se actualiza $Q(s,a)$ con la regla de Q-learning
    ($\alpha = 0.1$, $\gamma = 0.9$).
  \item \textbf{Fin del episodio:} cuando no quedan vehículos por
    insertar ni circulando (\verb|simulation.getMinExpectedNumber|
    igual a cero) o al llegar a 6000 pasos, se cierra SUMO con
    \verb|traci.close|, se leen del \texttt{tripinfo} las métricas del
    episodio (Sección~\ref{sec:u3}) y se guarda la tabla Q en disco.
\end{enumerate}

\noindent\textbf{Qué es un episodio y qué cambia entre episodios.}
Un episodio es una simulación completa de una hora de tráfico (0 a
3600~s de demanda, más el tiempo necesario para que salgan los últimos
vehículos) sobre la misma red y la misma demanda. Dentro de la hora el
agente toma del orden de 600 a 800 decisiones por semáforo y actualiza
su tabla Q en cada una. Entre un episodio y el siguiente cambian
exactamente tres cosas:

\begin{itemize}
  \item la probabilidad de exploración $\varepsilon$, que arranca en 1.0
    y se multiplica por 0.9 al terminar cada episodio hasta un piso de
    0.05 (lo alcanza en el episodio 30; con los decaimientos más lentos
    usados en los escenarios grandes, 0.96 y 0.95, lo alcanza en los
    episodios 75 y 60);
  \item la semilla de SUMO, que vale la semilla base más el número de
    episodio;
  \item la tabla Q, que el agente hereda del episodio anterior.
\end{itemize}

La demanda es idéntica en todos los episodios: los flujos se definen
con \verb|vehsPerHour|, que inserta los vehículos equiespaciados, de
modo que la semilla solo altera el comportamiento de los conductores
(su factor de velocidad individual y la imperfección $\sigma$ del modelo
de seguimiento), no cuántos vehículos llegan ni cuándo. No se ajustan
$\alpha$, $\gamma$ ni los pesos de la recompensa entre episodios.

El número de episodios se fija por el schedule de $\varepsilon$: los
necesarios para que llegue al piso, más al menos diez de consolidación
en los que el agente casi solo explota lo aprendido. Con decaimiento
0.9 son 40 episodios; con 0.96, 120; con 0.95, 100. Para comprobar que
ese criterio basta se registran por episodio el número de estados
distintos visitados y el mayor cambio de un valor Q, y cada cinco
episodios se evalúa la política congelada ($\varepsilon = 0$) sobre una
copia de la tabla Q, con una semilla fija ajena al entrenamiento y sin
alterar el generador aleatorio del entrenamiento. Las métricas de los
episodios de entrenamiento describen el proceso de aprendizaje, con
exploración y con la tabla Q cambiando dentro de la hora; el resultado
del agente se mide aparte, con la política congelada y varias semillas,
según el plan de pruebas de la Sección~\ref{sec:plan}.

El núcleo del ciclo, simplificado a partir de la función
\texttt{correr\_episodio} de \texttt{rl\_semaforo.py} (Anexo~B), es el
siguiente:

\begin{verbatim}
traci.start(["sumo", "-c", CFG, "--seed", str(seed),
             "--tripinfo-output", tripinfo])
verdes, ambar, dur_ambar, verde_min = fases_verdes()  # fases del programa
verde, t_verde, pasos = 0, 0, 0
poner_fase(verdes[verde], 100000)     # el agente decide cuando cambiar
s = estado(verde)
while traci.simulation.getMinExpectedNumber() > 0 and pasos < 6000:
    if modo == "train" and random.random() < eps:
        a = random.randint(0, 1)      # exploracion
    else:
        a = mejor_accion(Q, s)        # explotacion (desempate al azar)
    cambiar = (a == 1 and t_verde >= verde_min[verdes[verde]]) \
              or t_verde >= VERDE_MAX
    if cambiar:
        f = verdes[verde]
        poner_fase(ambar[f], dur_ambar[f])
        avanzar(dur_ambar[f])         # ambar de 3 s
        verde = (verde + 1) % len(verdes)
        poner_fase(verdes[verde], 100000)
        t_verde = 0
    avanzar(PASO_CONTROL)             # 5 s de simulacion
    t_verde += PASO_CONTROL
    r = recompensa()                  # -(w1 * colas + w2 * esperas)
    s2 = estado(verde)
    if modo == "train":
        q = q_vals(Q, s)
        q[a] += ALPHA * (r + GAMMA * max(q_vals(Q, s2)) - q[a])
    s = s2
traci.close()
\end{verbatim}

El mismo bucle sirve para los controles de referencia (tiempo fijo y
actuado): en esos modos el semáforo corre su propio programa y el
código solo mide, con la misma recompensa y las mismas métricas. El
script completo, con los modos \texttt{baseline}, \texttt{actuado},
\texttt{train} y \texttt{demo}, se incluye en los Anexos. El
entrenamiento puede adaptarse a algoritmos más avanzados sustituyendo
la tabla Q por una red neuronal (DQN) o por una política parametrizada
(PPO); el bucle de interacción con SUMO no cambia.

\subsection*{Algoritmos recomendados}

Diversas técnicas de RL han sido aplicadas al control de
semáforos. Los algoritmos tabulares, como \emph{Q-learning} o
\emph{SARSA}, son adecuados para intersecciones sencillas con un
número limitado de estados. Para escenarios de mayor complejidad se
emplean métodos de \emph{Deep Reinforcement Learning} como \emph{Deep
Q-Network} (DQN), \emph{Proximal Policy Optimization} (PPO) o
\emph{Deep Deterministic Policy Gradients} (DDPG), que utilizan redes
neuronales para aproximar la política y la función de valor. Estos
algoritmos suelen implementarse con bibliotecas de aprendizaje
automático (TensorFlow o PyTorch) y pueden integrarse con SUMO a
través de entornos compatibles con Gym. La elección del algoritmo
dependerá del tamaño del estado, la necesidad de cooperación entre
intersecciones y los recursos computacionales disponibles.

\subsection*{Actualización de los valores Q y ejemplo de pseudocódigo}

En el caso de algoritmos tabulares, como Q-learning, la política del agente se
representa mediante una \emph{tabla Q} en la que cada entrada $Q(s,a)$
almacena la expectativa de recompensa acumulada si se aplica la acción $a$ en
el estado $s$ y se sigue la mejor política a partir de entonces.
El proceso de aprendizaje consiste en actualizar iterativamente estos valores
con información obtenida de las simulaciones.  La regla de actualización de
Q-learning se basa en el aprendizaje por diferencias temporales y se expresa
de la siguiente manera:

\[
  Q(s,a) \leftarrow Q(s,a) + \alpha \bigl[r + \gamma\,\max_{a'} Q(s',a') - Q(s,a)\bigr],
\]
donde $\alpha$ es la tasa de aprendizaje, $\gamma$ el factor de descuento,
$r$ la recompensa inmediata obtenida al ejecutar la acción $a$ en el estado $s$,
$s'$ el estado siguiente y $a'$ las posibles acciones en $s'$.  Esta
ecuación desplaza el valor Q actual en dirección del retorno observado
($r + \gamma \max_{a'} Q(s',a')$).

Para algoritmos \emph{on--policy}, como SARSA, la actualización se realiza con la
siguiente fórmula:

\[
  Q(s,a) \leftarrow Q(s,a) + \alpha\bigl[r + \gamma\,Q(s',a') - Q(s,a)\bigr],
\]

donde $a'$ es la acción efectivamente seleccionada en el estado siguiente $s'$
y $\alpha,\gamma$ tienen el mismo significado que en Q-learning.  Al emplear
el valor $Q(s',a')$ asociado a la política vigente en vez del máximo de todas
las acciones, SARSA mantiene la coherencia entre exploración y actualización
y suele conducir a comportamientos más conservadores.

\subsection*{Diagrama de flujo del módulo U2}

La Figura~\ref{fig:flow-u2} sintetiza el ciclo de entrenamiento del módulo U2.
A partir del escenario generado en U1, el agente de RL interactúa con SUMO a
través de TraCI, observa el estado del tráfico, selecciona acciones sobre el
semáforo, actualiza sus valores $Q(s,a)$ (o los parámetros de la red neuronal)
y registra métricas de desempeño por episodio.
\begin{figure}[H]
    \centering
    \includegraphics[width=0.70\textwidth]{U2.png}
    \caption{Diagrama de flujo del módulo U2: entrenamiento del agente de
    Aprendizaje por Refuerzo para el control semafórico.}
    \label{fig:flow-u2}
\end{figure}

%----------------------------------------------------------------------
\section{Módulo U3: medición del desempeño}
\label{sec:u3}

Todas las métricas se calculan con las mismas salidas de SUMO para los
tres controles comparados (tiempo fijo, actuado y agente RL), de modo
que la comparación no dependa de cómo decide cada uno.

\begin{itemize}
  \item \textbf{Por vehículo}, del archivo \texttt{tripinfo} que SUMO
    escribe al terminar cada viaje: el retraso (\texttt{timeLoss}), la
    espera detenida (\texttt{waitingTime}), el número de paradas
    (\texttt{waitingCount}, veces que la velocidad bajó de 0.1~m/s) y
    el instante de llegada. El \textit{throughput} se define como el
    número de vehículos que completan su viaje antes de los 3600~s. Con
    el dispositivo de emisiones activado
    (\verb|--device.emissions.probability 1|) cada viaje incluye además
    sus emisiones acumuladas; se reportan CO$_2$ en gramos y NO$_x$ en
    miligramos por vehículo, con la clase de emisión por defecto de SUMO
    (HBEFA3, turismo a gasolina Euro~4) \citep{SUMO_Emissions}. Activar
    el dispositivo no altera la dinámica ni la secuencia aleatoria de la
    simulación: las métricas con y sin él coinciden.
  \item \textbf{Por segundo}, vía TraCI: el número de vehículos
    detenidos en las aproximaciones del cruce, del que salen la cola
    promedio y la cola máxima del episodio, y una serie por minuto para
    observar la respuesta de cada control al cambio de demanda.
  \item \textbf{Por decisión}: la recompensa de la
    ecuación~\ref{eq:recompensa} se evalúa cada 5~s también bajo los
    controles de referencia, para comparar lo que el agente optimiza
    con lo que consiguen el fijo y el actuado. Se reporta por decisión y
    no acumulada, porque el agente no muestrea durante los ámbares y
    acumula menos decisiones por hora. Se cuentan además los cambios de
    fase (entradas a una fase verde) de cada control.
\end{itemize}

Dos precisiones importan para leer los resultados. El retraso de SUMO
incluye tiempo perdido que no depende del semáforo (arranque desde
velocidad cero, giros a menor velocidad, vacilaciones del conductor),
así que un cambio porcentual sobre \texttt{timeLoss} subestima el
efecto sobre el retraso atribuible al control. Y la espera detenida es
la señal que optimiza la recompensa, pero no es el retraso: un control
que fragmenta las colas en paradas más cortas puede bajar la espera y
subir las paradas sin mover el retraso. Por eso el retraso es la
métrica primaria y la espera, las colas, las paradas, el
\textit{throughput} y las emisiones son secundarias.

%----------------------------------------------------------------------
\section{Plan de pruebas}
\label{sec:plan}

\subsection*{Casos, factores y controles}

Cada caso de prueba es un escenario de la Tabla~\ref{tab:escenarios}
sobre el que se comparan tres controles con la misma demanda y las
mismas semillas (Tabla~\ref{tab:plan}):

\begin{itemize}
  \item \textbf{Tiempo fijo tuneado:} el programa estático de
    \texttt{netconvert}, con el verde (y el desfase entre cruces)
    ajustado por barrido y comprobado por retraso promedio sobre varias
    semillas (Sección~\ref{sec:escenarios}).
  \item \textbf{Actuado:} el programa \verb|type="actuated"| nativo de
    SUMO con las mismas fases del fijo y con verde mínimo y máximo
    iguales a los del agente (10~s, 5~s en giros protegidos, 60~s); SUMO
    crea automáticamente los detectores. Cada cruce actúa por su cuenta.
  \item \textbf{Q-learning:} el agente de la Sección anterior con la
    política congelada ($\varepsilon = 0$) al terminar el entrenamiento.
\end{itemize}

Los factores variados son el escenario (cuatro niveles) y el control
(tres niveles). Se mantienen constantes la demanda (un nivel por
escenario, con la punta invertida a mitad de hora), el modelo de
llegadas (equiespaciadas) y la formulación e hiperparámetros del
agente. Quedan previstos y no ejecutados en esta fase: la demanda
escalada (0.8 y 1.2 veces la nominal), las llegadas aleatorias
(\verb|period="exp(...)"|), los algoritmos SARSA, DQN y PPO, la
coordinación del módulo U4 y el cruce real de Lima.

\begin{table}[H]
\centering
\footnotesize
\begin{tabular}{llcccl}
\toprule
Caso & Escenario & Agentes & \makecell{Episodios\\(decaimiento de $\varepsilon$)} & \makecell{Semillas de\\entrenamiento} & \makecell{Semillas de\\evaluación} \\
\midrule
P1 & cruce    & 1 & 40 (0.90)  & 43 a 82  & 1001 a 1010 \\
P2 & cruce2   & 1 & 120 (0.96) & 43 a 162 & 1001 a 1010 \\
P3 & corredor & 2 & 40 (0.90)  & 43 a 82  & 1001 a 1010 \\
P4 & red      & 3 & 100 (0.95) & 43 a 142 & 1001 a 1010 \\
P5 & cruce real de Lima & pendiente & -- & -- & -- \\
\bottomrule
\end{tabular}
\caption{Casos de prueba. En todos se comparan tiempo fijo tuneado, actuado y Q-learning; la métrica primaria es el retraso por vehículo y la meta, una reducción del 25\,\% frente al fijo.}
\label{tab:plan}
\end{table}

\subsection*{Réplicas y análisis estadístico}

El entrenamiento es una corrida por caso con semilla base 42 (el
episodio $k$ usa la semilla $42 + k$). Es determinista: repetir el
entrenamiento reproduce la misma tabla Q y las mismas métricas por
episodio. La evaluación se hace con la política congelada sobre diez
semillas (1001 a 1010) que no se usaron en el entrenamiento, las mismas
para los tres controles. La unidad experimental es la corrida de
evaluación y no el episodio de entrenamiento: dentro de un episodio la
tabla Q cambia y los episodios sucesivos no son independientes.

Por métrica y control se reportan la media, la desviación estándar y
el intervalo de confianza al 95\,\% (t de Student). Para cada control
frente al fijo se reportan el cambio porcentual de la media con su
intervalo al 95\,\% por \textit{bootstrap} pareado por semilla, la
prueba de Wilcoxon de rangos con signo (pareada por semilla), la prueba
de Mann--Whitney exacta y el delta de Cliff como tamaño de efecto. Se
declara una mejora cuando el intervalo del cambio excluye el cero y la
prueba pareada da $p < 0.05$, y se contrasta con las metas de la
Introducción.

\subsection*{Alcance de la medición}

Con la demanda actual todos los vehículos completan su viaje bajo los
tres controles y solo los que salen al final de la hora (entre 2 y
3\,\%) llegan después de los 3600~s, así que el \textit{throughput} no
discrimina entre controles y la meta de $+10\,\%$ no es evaluable en
estos escenarios; requiere un caso con demanda por encima de la
capacidad.
Las llegadas equiespaciadas hacen que la variabilidad entre semillas
provenga solo del comportamiento de los conductores, por lo que los
intervalos de confianza son estrechos; con llegadas aleatorias los
controles de referencia empeoran y los resultados no son comparables
con los de este plan, de modo que ese nivel se trata como un caso
aparte.

\subsection*{Reproducibilidad}

Cada corrida deja su fila en \texttt{resultados\_<escenario>.csv}
(entrenamiento y referencias) o en \texttt{evaluacion\_<escenario>.csv}
(evaluación por semilla), con la semilla y la configuración usadas; las
tablas Q entrenadas se guardan como JSON y las figuras se regeneran con
\texttt{graficar.py}. Entorno: SUMO 1.25.0 y Python 3.14 con
\texttt{traci}, \texttt{sumolib}, \texttt{numpy}, \texttt{scipy} y
\texttt{matplotlib}, sobre Windows 10 sin GPU; un episodio de una hora
simulada tarda unos pocos segundos y un entrenamiento completo, unos
minutos.
